\section{Scaling to Physical Mass Gap: The Program}
\label{sec:spectral-bound-program}

\textbf{Status: [PROGRAM] / [RESEARCH ROADMAP]}

Even if a mass gap $\Delta(\beta)$ is proven for all $\beta$, one must show that it scales correctly to a finite physical mass in the continuum limit. This requires relating the gap to a physical observable like the string tension $\sigma(\beta)$.

\subsection{Rigorous Foundation: Spectral Representation}

\begin{theorem}[Spectral Decomposition of Wilson Loop]
\label{thm:spectral-wilson-rigorous}
On a finite spatial lattice of size $L$, the Wilson loop has the spectral representation:
\[
\langle W_{R \times T} \rangle_L = \sum_{n=0}^\infty |\langle \Omega_L | \Phi_R | n \rangle|^2 e^{-E_n^{(L)} T}
\]
where $E_n^{(L)}$ are the eigenvalues of the Hamiltonian.
\end{theorem}

\subsection{The Target Inequality: Giles-Teper Bound}

To ensure the mass gap survives the continuum limit, we aim to prove:

\begin{conjecture}[Dimensionless Spectral Bound]
\label{conj:giles-teper}
There exists a universal constant $c > 0$ such that for all $\beta$ sufficiently large:
\[
\Delta(\beta) \ge c \sqrt{\sigma(\beta)}
\]
where $\sigma(\beta)$ is the string tension defined by the area law decay of large Wilson loops: $\langle W_{R \times T} \rangle \sim e^{-\sigma R T}$.
\end{conjecture}

\subsection{Implication for the Continuum Limit}

If Conjecture \ref{conj:giles-teper} holds, and if the string tension is physical (i.e., $\sigma_{\text{phys}} = \lim a^{-2} \sigma(\beta)$ is finite and non-zero), then:
\[
m_{\text{phys}} = \lim_{a \to 0} a^{-1} \Delta(\beta) \ge c \lim_{a \to 0} \sqrt{a^{-2} \sigma(\beta)} = c \sqrt{\sigma_{\text{phys}}} > 0
\]
This would prove the existence of a positive mass gap in the physical theory.

\subsection{Proposed Method}

The strategy to prove Conjecture \ref{conj:giles-teper} involves:
\begin{itemize}
    \item \textbf{Variational Bounds}: Constructing trial states for the Hamiltonian that approximate the "flux tube" state.
    \item \textbf{Spectral Analysis}: Analyzing the spectrum of the transfer matrix in the sector with non-trivial flux.
\end{itemize}
This connects the confinement property (area law) directly to the spectral gap.
